% GENERATED FILE. Edit canonical lessons, then run npm run generate:books.
% canonical-source-hash: fnv1a64:5fac2b169613f170
% canonical-lessons: ML-C01-namaskaram, ML-W01-na-ma-trace, ML-W01-na-ma-guided-copy, ML-W01-na-ma-delayed-copy, ML-W01-na-ma-dictation, ML-W01-sa-chandrakkala-ka, ML-W01-aa-ra-anusvaram, ML-W01-namaskaram-read, ML-W01-namaskaram-dictation, ML-C01-nandi, ML-C01-athe, ML-S133-letter-a, ML-C01-illa, ML-C01-sari, ML-C01-practice

\chapter{Greetings}
\label{ch:greetings}
\hlchaptermodality{1}

\begin{chapteropening}
\textbf{By the end of this chapter:} \emph{I can greet someone in Malayalam, read and write namaskāram from sound, recognize the reply, and choose an appropriate response.}

Greet and respond, read namaskāram aloud, then write it from dictation and check the three practised groups.
\end{chapteropening}

\section[namaskāram]{Namaskaram — first hear the greeting}

\label{lesson:ML-C01-namaskaram}



\begin{quote}
Imagine meeting someone with your palms together and a small bow. First, only listen: \emph{namaskāram}.
\end{quote}

\subsection*{The exchange}
\emph{Namaskāram} is a respectful Malayalam greeting. It can open a conversation, and it can also close one. Hear its five beats slowly: \emph{na-ma-s-kā-ram}. Answer with the same greeting once.

The Malayalam spelling is deliberately absent from this first pass. The next tiny lessons give your eye and hand no more than three reusable shapes at a time. Only after every shape is familiar will the book show the whole written greeting.

\begin{cousinweb}
The greeting comes from Sanskrit \emph{namas}, “a bow or homage,” plus \emph{kāra}, “a making,” from the root \emph{kṛ}, “to do.” Its memory hook is therefore \textbf{the making of a bow}. Malayalam is Tamil's closest sister, but this formal greeting is part of Malayalam's large Sanskrit layer.

The same idea, five ways:

\noindent
\begin{tabularx}{\linewidth}{@{}>{\raggedright\arraybackslash}X>{\raggedright\arraybackslash}X>{\raggedright\arraybackslash}X@{}}
\toprule
\textbf{Language} & \textbf{\textquotedblleft{}Hello\textquotedblright{}} & \textbf{Source} \\
\midrule
\textbf{Malayalam} & \emph{namaskāram} & Sanskrit \\
Telugu & \emph{namaskāram} & Sanskrit \\
Kannada & \emph{namaskāra} & Sanskrit \\
Hindi & \emph{namaste} & Sanskrit \\
Tamil & \emph{vaṇakkam} & \textbf{native Dravidian} \\
\bottomrule
\end{tabularx}

Even Tamil's closest sister took the Sanskrit greeting; Tamil alone kept its own. Hold that beside the next lesson, where the word for \textquotedblleft{}thanks\textquotedblright{} splits the family the other way round.
\end{cousinweb}

\subsection*{Your turn}
Say these aloud:

\begin{itemize}
\raggedright
  \item \emph{namaskāram} once as a greeting
  \item its five beats — \emph{na-ma-s-kā-ram}
  \item its literal memory hook — “the making of a bow”
\end{itemize}

\begin{tcolorbox}[breakable,colback=teal!4,colframe=teal!35!black,title={Before you move on}]
What can you say when you meet someone respectfully? (\emph{\textbf{Namaskāram.}}) What gesture and literal history help you remember it? (\textbf{A small bow; “the making of a bow.”}) Do you need to decode its Malayalam spelling yet? (\textbf{No.} The script starts next with only two shapes.)
\end{tcolorbox}
\section[na ma]{\ml{ന} · \ml{മ} — two sounds you already heard}

\label{lesson:ML-W01-na-ma-trace}



\begin{quote}
Say \emph{na-ma} — the first two beats of \emph{namaskāram}. Your hand meets only those two familiar sounds today.
\end{quote}

\subsection*{Script}
Malayalam is an abugida: a bare consonant already carries a short \emph{a}.

\begin{itemize}
\raggedright
  \item \textbf{\ml{ന}} is \emph{na}.
  \item \textbf{\ml{മ}} is \emph{ma}.
\end{itemize}

Put them together: \textbf{\ml{നമ}}, \emph{na-ma}. Keep both models visible.

\subsection*{Writing — observe and trace}
Finger-trace \textbf{\ml{ന}} once, then \textbf{\ml{മ}} once. Trace \textbf{\ml{നമ}} once from left to right. Say \emph{na-ma} after your finger stops. This is observation with a visible model, not a memory test.

\begin{tcolorbox}[breakable,colback=teal!4,colframe=teal!35!black,title={Before you move on}]
Point to \emph{na}. (\textbf{\ml{ന}}.) Point to \emph{ma}. (\textbf{\ml{മ}}.) What vowel did you write after either consonant? (\textbf{None} — the short \emph{a} is inherent.)
\end{tcolorbox}
\section[na-ma]{\ml{നമ} — one guided copy}

\label{lesson:ML-W01-na-ma-guided-copy}



\begin{quote}
Keep \textbf{\ml{നമ}} visible. Point left to right and say \emph{na-ma}.
\end{quote}

\subsection*{Writing — copy with help}
Copy \textbf{\ml{നമ}} once directly below the visible model. Look back as often as you need. Check \textbf{\ml{ന}} first and \textbf{\ml{മ}} second; repair only one place, then stop.

\begin{tcolorbox}[breakable,colback=teal!4,colframe=teal!35!black,title={Before you move on}]
Read your copy. Which longer greeting begins with these two beats? (\emph{\textbf{Namaskāram.}}) No new shape was added here: the work was moving from a finger trace to one supported copy.
\end{tcolorbox}
\section[na-ma]{\ml{നമ} — one delayed copy}

\label{lesson:ML-W01-na-ma-delayed-copy}



\begin{quote}
Look at \textbf{\ml{നമ}} for five seconds. Say \emph{na-ma}, then cover the model.
\end{quote}

\subsection*{Writing — retrieve after a delay}
Write the two shapes once while the model is covered. If one shape will not come, peek briefly, cover it again, and finish. This is retrieval practice, not a speed test.

\begin{tcolorbox}[breakable,colback=teal!4,colframe=teal!35!black,title={Before you move on}]
Uncover \textbf{\ml{നമ}} and compare. Repair one shape if needed. What changed from the last lesson? (\textbf{The model was hidden before writing.})
\end{tcolorbox}
\section[na-ma]{Two heard beats — write them}

\label{lesson:ML-W01-na-ma-dictation}



\begin{quote}
Cover the answer lower on this page. No new Malayalam shape is hiding here.
\end{quote}

\subsection*{Writing — short dictation}
\emph{Hear:} \emph{na-ma}

Write the two-beat unit from sound alone. Now uncover and compare: \textbf{\ml{നമ}}. Check \textbf{\ml{ന}} first and \textbf{\ml{മ}} second. Repair one place if needed.

\begin{tcolorbox}[breakable,colback=teal!4,colframe=teal!35!black,title={Before you move on}]
You have moved the same small unit through trace, visible copy, delayed copy, and dictation. Say the complete greeting it begins. (\emph{\textbf{Namaskāram.}})
\end{tcolorbox}
\section[sa, chandrakkala, ka]{\ml{സ} · \ml{്} · \ml{ക} — build ska without a leap}

\label{lesson:ML-W01-sa-chandrakkala-ka}



\begin{quote}
Write \textbf{\ml{നമ}} from the heard cue \emph{na-ma}. Now set it aside. This lesson adds exactly three shapes.
\end{quote}

\subsection*{Script}
\begin{itemize}
\raggedright
  \item \textbf{\ml{സ}} is \emph{sa}, with its inherent \emph{a}.
  \item \textbf{\ml{്}} is the \emph{chandrakkala}. Attach it to remove that inherent vowel: \textbf{\ml{സ്}}.
  \item \textbf{\ml{ക}} is \emph{ka}.
\end{itemize}

Join the vowel-less first consonant to the next one: \textbf{\ml{സ്ക}}, \emph{ska}. The cluster is not a fourth shape; it is the three pieces you just met working together.

\subsection*{Writing — observe and trace three pieces}
Keep the models visible. Trace \textbf{\ml{സ}}, then \textbf{\ml{്}}, then \textbf{\ml{ക}} once each. Point through \textbf{\ml{സ്ക}} and say \emph{ska}. Stop after one slow pass.

\begin{tcolorbox}[breakable,colback=teal!4,colframe=teal!35!black,title={Before you move on}]
What does \textbf{\ml{്}} remove? (\textbf{The inherent vowel.}) Read \textbf{\ml{സ}} and then \textbf{\ml{സ്}}. (\emph{\textbf{sa; s.}}) What does \textbf{\ml{സ്ക}} say? (\emph{\textbf{ska.}})
\end{tcolorbox}
\section[long ā, ra, final m]{\ml{ാ} · \ml{ര} · \ml{ം} — finish the sound inventory}

\label{lesson:ML-W01-aa-ra-anusvaram}



\begin{quote}
Read \textbf{\ml{സ്ക}} once: \emph{ska}. The greeting needs only three more written pieces.
\end{quote}

\subsection*{Script}
\begin{itemize}
\raggedright
  \item \textbf{\ml{ാ}} is the long-\emph{ā} vowel sign. Attach it to known \textbf{\ml{ക}}: \textbf{\ml{കാ}}, \emph{kā}.
  \item \textbf{\ml{ര}} is \emph{ra}.
  \item \textbf{\ml{ം}} is the \emph{anusvaram}, the final nasal \emph{m} here. With \textbf{\ml{ര}}: \textbf{\ml{രം}}, \emph{ram}.
\end{itemize}

These are the last new shapes in the greeting.

\subsection*{Writing — observe and trace the last three}
Trace \textbf{\ml{ാ}}, \textbf{\ml{ര}}, and \textbf{\ml{ം}} once each while the models remain visible. Then trace \textbf{\ml{കാ}} and \textbf{\ml{രം}} once. Say \emph{kā-ram} after your finger stops.

\begin{tcolorbox}[breakable,colback=teal!4,colframe=teal!35!black,title={Before you move on}]
What makes \textbf{\ml{ക}} into \textbf{\ml{കാ}}? (\textbf{The long-\emph{ā} sign, \ml{ാ}.}) What does the final circle \textbf{\ml{ം}} contribute here? (\textbf{A final nasal \emph{m}.}) How many new shapes remain before reading the whole greeting? (\textbf{None.})
\end{tcolorbox}
\section[namaskāram]{\ml{നമസ്കാരം} — every shape is already yours}

\label{lesson:ML-W01-namaskaram-read}



\begin{quote}
Say \emph{namaskāram}. Now the written greeting may appear: it contains no unfamiliar shape.
\end{quote}

\subsection*{Script — assemble and read}
\begin{quote}
\textbf{\ml{ന}} + \textbf{\ml{മ}} + \textbf{\ml{സ}} + \textbf{\ml{്}} + \textbf{\ml{ക}} + \textbf{\ml{ാ}} + \textbf{\ml{ര}} + \textbf{\ml{ം}} $\to$ \textbf{\ml{നമസ്കാരം}}
\end{quote}

Read the pieces in the groups you practised:

\begin{itemize}
\raggedright
  \item \textbf{\ml{നമ}} — \emph{na-ma}
  \item \textbf{\ml{സ്ക}} — \emph{ska}
  \item \textbf{\ml{ാരം}} — \emph{ā-ram}
\end{itemize}

Together: \textbf{\ml{നമസ്കാരം}}, \emph{namaskāram}. The conjunct \textbf{\ml{സ്ക}} is still just \textbf{\ml{സ}} with its vowel removed by \textbf{\ml{്}}, joined to \textbf{\ml{ക}}.

\subsection*{Writing — one visible whole-word copy}
Keep \textbf{\ml{നമസ്കാരം}} visible and copy it once in the three groups above. Look back whenever needed. Independent production waits for the next tiny lesson.

\begin{tcolorbox}[breakable,colback=teal!4,colframe=teal!35!black,title={Before you move on}]
Read \textbf{\ml{നമസ്കാരം}}. How many new shapes did this lesson add? (\textbf{None.}) What does the greeting literally evoke? (\textbf{The making of a bow.})
\end{tcolorbox}
\section[namaskāram]{One heard greeting — write the whole}

\label{lesson:ML-W01-namaskaram-dictation}



\begin{quote}
Cover the written answer below. You have already traced every shape, carried \textbf{\ml{നമ}} through dictation, and copied the whole greeting once.
\end{quote}

\subsection*{Writing — whole-word dictation}
\emph{Hear:} \emph{namaskāram} — “greetings”

Write the greeting from sound alone. Use the three remembered groups: \emph{na-ma}, \emph{ska}, \emph{ā-ram}. Then uncover and compare: \textbf{\ml{നമസ്കാരം}}. Repair one group if needed; do not rewrite the whole word.

\begin{tcolorbox}[breakable,colback=teal!4,colframe=teal!35!black,title={Before you move on}]
Read your greeting and say what it does. (\textbf{It respectfully opens or closes a conversation.}) This is independent production of a fully taught word, not a preview of untaught shapes.
\end{tcolorbox}
\section[nandi]{\ml{നന്ദി} (nandi) — \textquotedblleft{}thank you,\textquotedblright{} and the Tamil kinship made visible}

\label{lesson:ML-C01-nandi}



\begin{quote}
After a Sanskrit greeting, the native core shows itself: Malayalam's \textquotedblleft{}thank you\textquotedblright{} is \textbf{not} the Sanskrit word its neighbours Kannada and Telugu use. It is \emph{nandi} — the very cousin of Tamil's \emph{naṉṟi}.
\end{quote}

\subsection*{The letters in this word}
\emph{(Skim if you read Malayalam.)}

\begin{itemize}
\raggedright
  \item \textbf{\ml{ന}} = \textquotedblleft{}na.\textquotedblright{}
  \item \textbf{\ml{ന്ദി}} = \textquotedblleft{}ndi\textquotedblright{} — a \textbf{conjunct}: \textbf{\ml{ന}} (na) loses its vowel and joins \textbf{\ml{ദ}} (da) into \textquotedblleft{}nd,\textquotedblright{} then the \textbf{\textquotedblleft{}i\textquotedblright{} vowel sign} gives \textquotedblleft{}ndi.\textquotedblright{}
\end{itemize}

Left to right: \textbf{\ml{ന} · \ml{ന്ദി}} = \emph{na-ndi} $\to$

\begin{quote}
\textbf{\ml{നന്ദി}} = \textbf{nandi}
\end{quote}

\begin{cousinweb}
\textbf{\ml{നന്ദി}} (\emph{nandi}) is \textbf{native Dravidian}, from the same root \textbf{nal / nan} (\textquotedblleft{}good\textquotedblright{}) that gives Tamil \textbf{naṉṟi} (\textquotedblleft{}goodness $\to$ gratitude\textquotedblright{}). Malayalam and Tamil split from one language about a thousand years ago, so their oldest, most everyday words are near-twins — and \textquotedblleft{}thanks\textquotedblright{} is a perfect example: Malayalam \emph{nandi}, Tamil \emph{naṉṟi}, one word lightly reshaped.

Its neighbours went the other way: Kannada and Telugu use the Sanskrit \emph{dhanyavāda}. So the map of \textquotedblleft{}thanks\textquotedblright{} traces the family's fault-line exactly — Tamil and Malayalam on the native side, Kannada and Telugu on the Sanskrit side.

The same idea, five ways:

\noindent
\begin{tabularx}{\linewidth}{@{}>{\raggedright\arraybackslash}X>{\raggedright\arraybackslash}X>{\raggedright\arraybackslash}X@{}}
\toprule
\textbf{Language} & \textbf{\textquotedblleft{}Thanks\textquotedblright{}} & \textbf{Note} \\
\midrule
\textbf{Malayalam} & \emph{nandi} (\ml{നന്ദി}) & \textbf{native} (root \emph{nal}, \textquotedblleft{}good\textquotedblright{}) \\
Tamil & \emph{naṉṟi} (\ta{நன்றி}) & the \textbf{same} native root \\
Kannada & \emph{dhanyavāda} (\ka{ಧನ್ಯವಾದ}) & Sanskrit \\
Telugu & \emph{dhanyavādamulu} (\te{ధన్యవాదములు}) & Sanskrit \\
Hindi & \emph{dhanyavād} / \emph{shukriyā} & Sanskrit / Perso-Arabic \\
\bottomrule
\end{tabularx}
\end{cousinweb}

\begin{culture}
Like \emph{naṉṟi} in Tamil, \emph{nandi} is a touch formal — among close family it can feel oddly stiff, as if you were thanking a stranger; warmth is usually assumed and shown in tone. The explicit word is for when you want to \emph{mark} the gratitude. (In casual speech you'll also hear the English \emph{thanks}.)
\end{culture}

\subsection*{Your turn}
Say these aloud:

\begin{itemize}
\raggedright
  \item na · ndi $\to$ \textquotedblleft{}nandi\textquotedblright{}
  \item the pair across the split — Malayalam \emph{nandi}, Tamil \emph{naṉṟi}
  \item how \ml{ന്ദ} joins two consonants into one \textquotedblleft{}nd\textquotedblright{} block
\end{itemize}

\begin{tcolorbox}[breakable,colback=teal!4,colframe=teal!35!black,title={Before you move on}]
Read \textbf{\ml{നന്ദി}}. Is it native or a Sanskrit loan, and which language has the twin word? (Native — root \emph{nal}, \textquotedblleft{}good\textquotedblright{}; Tamil's \emph{naṉṟi}.) Which two Dravidian languages use the \emph{Sanskrit} word for thanks instead? (Kannada, Telugu.)
\end{tcolorbox}
\section[athe]{\ml{അതെ} (athe) — \textquotedblleft{}yes,\textquotedblright{} literally \textquotedblleft{}that {[}is so{]}\textquotedblright{}}

\label{lesson:ML-C01-athe}



\begin{quote}
The plain \textquotedblleft{}yes\textquotedblright{} — and a small window into how Malayalam thinks, because the word is really the pointing-word \textquotedblleft{}that.\textquotedblright{}
\end{quote}

\subsection*{The letters in this word}
\emph{(Skim if you read Malayalam.)}

\begin{itemize}
\raggedright
  \item \textbf{\ml{അ}} = \textbf{\textquotedblleft{}a\textquotedblright{}} — an \textbf{independent vowel} (the word starts with a vowel, so it gets a full letter).
  \item \textbf{\ml{തെ}} = \textquotedblleft{}the\textquotedblright{} — \textbf{\ml{ത}} (tha, a soft \emph{th} as in \emph{this}) with the \textbf{\textquotedblleft{}e\textquotedblright{} vowel sign} written \emph{in front} of the consonant.
\end{itemize}

Left to right: \textbf{\ml{അ} · \ml{തെ}} = \emph{a-the} $\to$

\begin{quote}
\textbf{\ml{അതെ}} = \textbf{athe} = \textquotedblleft{}yes.\textquotedblright{}
\end{quote}

\begin{cousinweb}
\textbf{\ml{അതെ}} (\emph{athe}) is native Malayalam, and it grows from \textbf{\ml{അത്}} (\emph{athu}), \textquotedblleft{}that.\textquotedblright{} To say \textquotedblleft{}yes,\textquotedblright{} Malayalam essentially says \textbf{\textquotedblleft{}that {[}is so{]}\textquotedblright{}} — you affirm by pointing at the truth of what was said. Its natural partner is the negative \emph{alla}, \textquotedblleft{}{[}it is{]} not that\textquotedblright{} — the next lesson.
\end{cousinweb}

\begin{grammarlens}[title={Grammar Lens: yes and no as \textquotedblleft{}that is\textquotedblright{} / \textquotedblleft{}not that\textquotedblright{}}]
Malayalam's yes/no pair are demonstratives in disguise: \emph{athe} (\textquotedblleft{}that {[}is so{]}\textquotedblright{}) and \emph{alla} (\textquotedblleft{}not that\textquotedblright{}). Like its sisters, Malayalam also readily answers by \textbf{echoing the verb} of the question (\emph{vannu}, \textquotedblleft{}{[}I{]} came,\textquotedblright{} for \textquotedblleft{}did you come?\textquotedblright{}). Bank \emph{athe} as the simple word, but hear the logic underneath — you agree by confirming \emph{that} is how things are.
\end{grammarlens}

\subsection*{Your turn}
\begin{itemize}
\raggedright
  \item \emph{Say it:} a · the $\to$ \textquotedblleft{}athe\textquotedblright{}
  \item \emph{Say it:} note the \textquotedblleft{}e\textquotedblright{} sign sits \emph{before} \ml{ത} but is read \emph{after} it
  \item \emph{Say it:} \textquotedblleft{}athe\textquotedblright{} — \textquotedblleft{}yes, that's so\textquotedblright{}
  \item \emph{Recall:} say \emph{namaskāram}
\end{itemize}

\begin{tcolorbox}[breakable,colback=teal!4,colframe=teal!35!black,title={Before you move on}]
Read \textbf{\ml{അതെ}}. What everyday word is it built from, and so what does \textquotedblleft{}yes\textquotedblright{} literally assert? (\emph{athu}, \textquotedblleft{}that\textquotedblright{} — so \textquotedblleft{}yes\textquotedblright{} = \textquotedblleft{}that {[}is so{]}\textquotedblright{}.) Where does the \textquotedblleft{}e\textquotedblright{} vowel sign sit, and where is it read? (Written before the consonant, read after.)
\end{tcolorbox}
\section[a]{\ml{അ} — one character, met inside words you already say}

\label{lesson:ML-S133-letter-a}



\begin{quote}
Before the new one: \ml{നമസ്കാരം} — can you still read it back?

One character this time. One only. It has been sitting in front of you on pages you have already read.
\end{quote}

\subsection*{Script you'll notice: \ml{അ}}
\textbf{\ml{അ}} — \emph{a}.

It is an \textbf{independent vowel} — the shape this vowel takes when a word \emph{begins} with it, rather than the sign it becomes inside a word.

You have already said a word that starts with it.

\begin{itemize}
\raggedright
  \item \textbf{\ml{അതെ}} \emph{athe} — yes (\textquotedblleft{}that {[}is so{]}\textquotedblright{})
\end{itemize}

Malayalam writes a word-initial vowel with a letter of its own. Inside a word this vowel would be a small sign hung on a consonant; at the front of a word there is no consonant to hang it on, so the vowel stands up as \textbf{\ml{അ}}.

\subsection*{Writing: \ml{അ} — copy what you see}
Put your pen on \ml{അ} and follow its line. Copy the shape you can see — slowly, and larger than it is printed.

\begin{quote}
This book does not yet tell you \textbf{where to start the character or which way to travel}. That is a real thing, taught with real variation from school to school, and it is not written down here until it can be written down with a source. Copying what is in front of you needs no such source, and it is how the shape gets into your hand in the meantime.
\end{quote}

\subsection*{Your turn}
\begin{itemize}
\raggedright
  \item \emph{Look:} at these words, and find \ml{അ} in the ones that have it
\end{itemize}

\begin{quote}
\ml{അതെ}  ·  \ml{നന്ദി}  ·  \ml{നമസ്കാരം}
\end{quote}

\begin{itemize}
\raggedright
  \item \emph{Trace it:} \ml{അ} three times, saying \emph{a} as you finish each one
  \item \emph{Look:} back at any page of this chapter and find \ml{അ} once more
\end{itemize}

\begin{tcolorbox}[breakable,colback=teal!4,colframe=teal!35!black,title={Before you move on}]
Which character is this — \ml{അ}? What sound does it carry? (\emph{\textbf{a}}.) Where in a word does a vowel need a full letter of its own? (At the start.)
\end{tcolorbox}
\section[illa]{\ml{ഇല്ല} (illa) — \textquotedblleft{}no,\textquotedblright{} the word Malayalam and Tamil share almost exactly}

\label{lesson:ML-C01-illa}



\begin{quote}
The partner of \emph{athe}, and more family evidence: this Malayalam \textquotedblleft{}no\textquotedblright{} is essentially the \textbf{same word} as Tamil's \emph{illai} and Kannada's \emph{illa} — the old Dravidian \emph{il-} runs right through.
\end{quote}

\subsection*{The letters in this word}
\emph{(Skim if you read Malayalam.)}

\begin{itemize}
\raggedright
  \item \textbf{\ml{ഇ}} = \textbf{\textquotedblleft{}i\textquotedblright{}} — an \textbf{independent vowel} (word-initial, so a full letter, like \textbf{\ml{അ}} in \emph{athe}).
  \item \textbf{\ml{ല്ല}} = \textquotedblleft{}lla\textquotedblright{} — a \textbf{conjunct}: \textbf{\ml{ല}} (la) loses its vowel and joins another \textbf{\ml{ല}} (la) into a doubled, held \textquotedblleft{}ll.\textquotedblright{}
\end{itemize}

Left to right: \textbf{\ml{ഇ} · \ml{ല്ല}} = \emph{i-lla} $\to$

\begin{quote}
\textbf{\ml{ഇല്ല}} = \textbf{illa}
\end{quote}

\begin{cousinweb}
\textbf{\ml{ഇല്ല}} (\emph{illa}) is built on the ancient Dravidian root \textbf{il}, \textquotedblleft{}not-being, absence\textquotedblright{} — so, exactly like Tamil \emph{illai}, it really means \textbf{\textquotedblleft{}{[}it{]} is not, {[}there{]} is not,\textquotedblright{}} and serves as the everyday \textquotedblleft{}no.\textquotedblright{} The positive partner will come much later, after you have enough sentence pieces to use the contrast.

The same idea, five ways:

\noindent
\begin{tabularx}{\linewidth}{@{}>{\raggedright\arraybackslash}X>{\raggedright\arraybackslash}X>{\raggedright\arraybackslash}X@{}}
\toprule
\textbf{Language} & \textbf{\textquotedblleft{}No / isn't\textquotedblright{}} & \textbf{Note} \\
\midrule
\textbf{Malayalam} & \emph{illa} (\ml{ഇല്ല}) & native, root \emph{il-} \\
Tamil & \emph{illai} (\ta{இல்லை}) & the \textbf{same} root \\
Kannada & \emph{illa} (\ka{ಇಲ್ಲ}) & the \textbf{same} root \\
Telugu & \emph{lēdu} (\te{లేదు}) & native, a \textbf{different} root \\
Hindi & \emph{nahīṁ} (\dv{नहीं}) & Indo-European (unrelated) \\
\bottomrule
\end{tabularx}
\end{cousinweb}

\begin{grammarlens}[title={Grammar Lens: Dravidian negates by \emph{being}}]
English inserts \emph{not} / \emph{no}. Malayalam — like Tamil and Kannada — can instead use a \textbf{negative verb of existence}. So \emph{illa} is both the plain \textquotedblleft{}no\textquotedblright{} and, later, a word that can make a whole sentence mean \textquotedblleft{}there is not.\textquotedblright{} Today, only the short answer is yours to use. This is a \textbf{shared Dravidian inheritance}, strongest across the Tamil–Malayalam pair: Tamil \emph{illai}, Malayalam \emph{illa}, Kannada \emph{illa} — one family, one way to say a thing is not. (Recall Telugu was the outlier here, using \emph{lēdu}.)
\end{grammarlens}

\subsection*{Your turn}
Say these aloud:

\begin{itemize}
\raggedright
  \item i · lla $\to$ \textquotedblleft{}illa\textquotedblright{}
  \item the pair — \emph{athe} (\textquotedblleft{}yes\textquotedblright{}) / \emph{illa} (\textquotedblleft{}no\textquotedblright{})
  \item \emph{illa} once as the short answer \textquotedblleft{}no\textquotedblright{}
\end{itemize}

\begin{tcolorbox}[breakable,colback=teal!4,colframe=teal!35!black,title={Before you move on}]
Read \textbf{\ml{ഇല്ല}}. Literally, is it closer to \textquotedblleft{}no\textquotedblright{} or \textquotedblleft{}is not\textquotedblright{}? (\textquotedblleft{}Is not / there is not.\textquotedblright{}) Which languages share this exact \emph{il-} root, and which Dravidian sister does \emph{not}? (Tamil and Kannada share it; Telugu is the outlier, with \emph{lēdu}.)
\end{tcolorbox}
\section[sari]{\ml{ശരി} (śari) — \textquotedblleft{}okay,\textquotedblright{} the whole-family word once more}

\label{lesson:ML-C01-sari}



\begin{quote}
The easy middle between \textquotedblleft{}yes\textquotedblright{} and \textquotedblleft{}no\textquotedblright{}: \textquotedblleft{}okay, alright, correct.\textquotedblright{} You have now met it in four languages — Tamil and Kannada \emph{sari}, Telugu \emph{sarē}, and here Malayalam \emph{śari}.
\end{quote}

\subsection*{The letters in this word}
\emph{(Skim if you read Malayalam.)}

\begin{itemize}
\raggedright
  \item \textbf{\ml{ശ}} = \textquotedblleft{}śa\textquotedblright{} — a \emph{sh}-sound (the palatal \emph{ś}). Malayalam, having borrowed so much Sanskrit, keeps the full set of Sanskrit sibilants (\emph{ś, ṣ, s}) as separate letters — more than Tamil, which makes do with one.
  \item \textbf{\ml{രി}} = \textquotedblleft{}ri\textquotedblright{} — \textbf{\ml{ര}} (ra) with the \textbf{\textquotedblleft{}i\textquotedblright{} vowel sign}.
\end{itemize}

Left to right: \textbf{\ml{ശ} · \ml{രി}} = \emph{śa-ri} $\to$

\begin{quote}
\textbf{\ml{ശരി}} = \textbf{śari} = \textquotedblleft{}okay / correct.\textquotedblright{}
\end{quote}

\begin{cousinweb}
\textbf{\ml{ശരി}} (\emph{śari}) is the pan-Dravidian \textbf{\textquotedblleft{}correct, right $\to$ okay, agreed,\textquotedblright{}} the verbal nod of Malayalam. It is the same family word carried by Tamil and Kannada as \emph{sari} and Telugu as \emph{sarē}; Malayalam simply spells its opening sound with \textbf{\ml{ശ}} (\emph{ś}). Doubled, \emph{śari śari}, it means \textquotedblleft{}yes yes, fine.\textquotedblright{}
\end{cousinweb}

\begin{grammarlens}[title={Grammar Lens: one word, four scripts}]
\emph{sari / sarē / śari} is the clearest single proof of how close these four languages are: essentially one word and one meaning, wearing four different scripts. As you learn to read each, catching a shared word change its clothes but not its sense is the fastest way to feel the whole Dravidian family at once — and the reason the four are worth seeing side by side.
\end{grammarlens}

\subsection*{Your turn}
Say these aloud:

\begin{itemize}
\raggedright
  \item śa · ri $\to$ \textquotedblleft{}śari\textquotedblright{}
  \item the everyday triple — \emph{athe} (yes) · \emph{illa} (no) · \emph{śari} (okay)
  \item the word across the family — \emph{sari} / \emph{sarē} / \emph{śari}
\end{itemize}

\begin{tcolorbox}[breakable,colback=teal!4,colframe=teal!35!black,title={Before you move on}]
Read \textbf{\ml{ശരി}}. What does it mean, and how do its Tamil, Kannada, and Telugu twins sound? (\textquotedblleft{}Okay, correct\textquotedblright{}; \emph{sari}, \emph{sari}, \emph{sarē}.) Why does Malayalam keep a separate \textbf{\ml{ശ}} (\emph{ś}) letter where Tamil uses one all-purpose sibilant? (It borrowed heavily from Sanskrit and kept Sanskrit's sibilant letters.)
\end{tcolorbox}
\section[Practice]{Chapter 1 — a five-word recap}

\label{lesson:ML-C01-practice}



\begin{quote}
Five words, and a first handful of Malayalam letters read straight off the page. Let's gather only what you have already learned and make it feel easy.
\end{quote}

\subsection*{The exchange — read them back}
Sound each out, left to right, before checking:

\noindent
\begin{tabularx}{\linewidth}{@{}>{\raggedright\arraybackslash}X>{\raggedright\arraybackslash}X>{\raggedright\arraybackslash}X>{\raggedright\arraybackslash}X@{}}
\toprule
\textbf{Read} & \textbf{} & \textbf{Meaning} & \textbf{Origin} \\
\midrule
\textbf{\ml{നമസ്കാരം}} & \emph{namaskāram} & hello / greetings & Sanskrit \\
\textbf{\ml{നന്ദി}} & \emph{nandi} & thank you & native (= Tamil \emph{naṉṟi}) \\
\textbf{\ml{അതെ}} & \emph{athe} & yes & native (\textquotedblleft{}that {[}is so{]}\textquotedblright{}) \\
\textbf{\ml{ഇല്ല}} & \emph{illa} & no / there isn't & native (= Tamil \emph{illai}) \\
\textbf{\ml{ശരി}} & \emph{śari} & okay & native (= \emph{sari}) \\
\bottomrule
\end{tabularx}

The letters that did it: consonants carry a built-in \textbf{\textquotedblleft{}a\textquotedblright{}}; the \textbf{chandrakkala} strips that vowel so consonants can \textbf{join into conjuncts} (\ml{സ്ക}, \ml{ന്ദ}, \ml{ല്ല}); a \textbf{vowel sign} swaps the vowel (\ml{ാ} ā, \ml{ി} i, \ml{െ} e — the \emph{e}-sign written \emph{before} its consonant); the \textbf{anusvāram} (\ml{ം}) adds a final nasal; and a word-initial vowel gets a \textbf{full letter} (\ml{അ}, \ml{ഇ}).

Notice how \textbf{Tamil-like} the native core is: four of the five everyday words (\emph{nandi, athe, illa, śari}) are shared with Tamil — because Malayalam \emph{is} Tamil's closest sister — while only the formal greeting reached for Sanskrit.

\begin{culture}
Malayalam has a warm departure phrase built around going and coming back. It belongs in the later chapter where you learn those actions one at a time. For now, repeating the greeting is enough when you part.
\end{culture}

\subsection*{Your turn}
Say these aloud:

\begin{itemize}
\raggedright
  \item all five — namaskāram · nandi · athe · illa · śari
  \item greet and thank — \textquotedblleft{}namaskāram!\textquotedblright{} … \textquotedblleft{}nandi.\textquotedblright{}
\end{itemize}

\subsection*{Writing checkpoint}
Hide the model. Hear \emph{namaskāram} once, write \textbf{\ml{നമസ്കാരം}}, then uncover the word in the table and check one group at a time: \textbf{\ml{നമ} · \ml{സ്ക} · \ml{ാരം}}. A mismatch is a signal to repeat the earlier two-minute shape lesson, not to push faster.

\subsection*{Script check — the letters of this chapter}
Cover the romanizations. Find each letter once in the Malayalam above, then say the word it sits in.

\begin{itemize}
\raggedright
  \item \textbf{\ml{അതെ}} — \textbf{\ml{അ}} \emph{a}
\end{itemize}

\begin{itemize}
\raggedright
  \item \emph{Look:} at each letter, then find it in the word beside it
\end{itemize}

\begin{tcolorbox}[breakable,colback=teal!4,colframe=teal!35!black,title={Before you move on}]
Read all five aloud. Which one is Sanskrit, and which four are native Dravidian shared with Tamil? (Sanskrit: \emph{namaskāram}. Native: \emph{nandi, athe, illa, śari}.) Which word can mean both \textquotedblleft{}okay\textquotedblright{} and a casual sign-off? (\emph{śari}.)

That completes the four Dravidian first chapters — Tamil, Kannada, Telugu, Malayalam. Next chapter starts introducing yourself, one short possession form at a time.
\end{tcolorbox}
